\documentclass[11pt,a4paper]{article}

\usepackage[margin=1.05in]{geometry}
\usepackage{amsmath,amssymb,amsthm}
\usepackage{mathtools}
\usepackage{stmaryrd}
\usepackage{microtype}
\setlength{\emergencystretch}{2.5em}
\usepackage{enumitem}
\usepackage{booktabs}
\usepackage{array}
\usepackage{tikz-cd}
\usepackage{xcolor}
\usepackage{listings}
\usepackage[colorlinks=true,linkcolor=blue!45!black,urlcolor=blue!60!black]{hyperref}

\newtheorem{theorem}{Theorem}[section]
\newtheorem{proposition}[theorem]{Proposition}
\newtheorem{lemma}[theorem]{Lemma}
\newtheorem{corollary}[theorem]{Corollary}
\newtheorem{conjecture}[theorem]{Conjecture}
\theoremstyle{definition}
\newtheorem{definition}[theorem]{Definition}
\newtheorem{example}[theorem]{Example}
\theoremstyle{remark}
\newtheorem{remark}[theorem]{Remark}
\newtheorem{warning}[theorem]{Warning}
\newtheorem{refuted}[theorem]{Refuted conjecture}
\newtheorem{question}[theorem]{Question}

\newcommand{\FF}{\mathbb{F}}
\newcommand{\CC}{\mathbb{C}}
\newcommand{\NN}{\mathbb{N}}
\newcommand{\ZZ}{\mathbb{Z}}
\newcommand{\cat}[1]{\mathbf{#1}}
\newcommand{\Set}{\cat{Set}}
\newcommand{\FinSet}{\cat{FinSet}}
\newcommand{\Vect}{\cat{Vect}}
\newcommand{\finVect}{\cat{finVect}}
\newcommand{\Cont}{\cat{Cont}}
\newcommand{\Rep}{\mathrm{Rep}}
\newcommand{\ext}[1]{\llbracket #1 \rrbracket}
\newcommand{\con}{\triangleleft}
\newcommand{\op}{\mathrm{op}}
\DeclareMathOperator{\Hom}{Hom}
\DeclareMathOperator{\Nat}{Nat}
\DeclareMathOperator{\GL}{GL}
\DeclareMathOperator{\im}{im}
\DeclareMathOperator{\End}{End}
\DeclareMathOperator{\Ind}{Ind}
\DeclareMathOperator{\Res}{Res}
\DeclareMathOperator{\Coind}{Coind}
\DeclareMathOperator{\Infl}{Infl}
\newcommand{\Repg}{\Rep_{k}(G)}
\newcommand{\qbin}[2]{\genfrac{[}{]}{0pt}{}{#1}{#2}_{q}}
\newcommand{\yon}{\mathsf{y}}

\definecolor{okgreen}{rgb}{0.0,0.45,0.15}
\definecolor{nored}{rgb}{0.65,0.0,0.0}
\newcommand{\OK}{\textcolor{okgreen}{\textbf{confirmed}}}
\newcommand{\NO}{\textcolor{nored}{\textbf{refuted}}}

\title{\bfseries Containers over a Finite Field\\[0.4em]
\large first definitions, two theorems, and one refuted conjecture}
\author{Robin Langer \and Claude}
\date{11 September 2026}

\begin{document}
\maketitle

\begin{abstract}
A container $S \con P$ (Abbott--Altenkirch--Ghani) is a set of shapes together
with a set of positions for each shape; its extension
$\ext{S \con P}X = \sum_{s} X^{P s}$ is a polynomial functor, and every strictly
positive type arises this way. We observe that a container is a piece of
\emph{$\FF_1$-linear algebra}: positions form a set, and the symmetry group of $n$
positions is $S_n = \GL_n(\FF_1)$. This document deforms that structure to a
general finite field.

The naive internalisation fails --- $\Vect$ is additive, so cartesian-polynomial
functors on it collapse to linear ones --- and the correct move is to observe that
a container is a \emph{set-indexed coproduct of corepresentables} and to change
only the base category. The resulting notion, the \emph{$q$-container}, has
positions that are $\FF_q$-vector spaces and reindexing that is linear.

We prove a representation theorem (Theorem~\ref{thm:A}), prove that linearising
the shapes destroys the theory completely (Theorem~\ref{thm:B}), identify the
derivative with the Jackson $q$-derivative (Theorem~\ref{thm:D}), and
\emph{refute} our own conjecture that $q$-containers carry two derivatives
(\S\ref{sec:refuted}) --- the refutation locating the interesting deformation not
in $q$ but in the self-duality of the base category. All results are
computationally verified by exact enumeration over $\FF_2$ and $\FF_3$;
\S\ref{sec:status} records exactly which claims are proved, which are verified,
and which remain open.

\medskip
\noindent\emph{Companion document:} \texttt{finite-field-duality.pdf}, which
supplies the background on forms over $\FF_q$, the topos/abelian dichotomy, and
the $q \to 1$ limit; references of the form ``FFD \S$n$'' point there.
\end{abstract}

\tableofcontents

\section{Containers, and the observation that starts this}

\begin{definition}[Container]
A \emph{container} is a pair $S \con P$ where $S$ is a set of \emph{shapes} and
$P : S \to \Set$ assigns to each shape a set of \emph{positions}. Its
\emph{extension} is the functor $\Set \to \Set$
\[
  \ext{S \con P} X \;=\; \sum_{s : S} \bigl(P s \to X\bigr)
  \;=\; \sum_{s:S} X^{P s},
\]
``choose a shape, then label every position with an element of $X$''.
\end{definition}

\begin{definition}[Container morphism]\label{def:contmor}
A morphism $(S \con P) \to (T \con Q)$ is a pair
\[
  u : S \to T, \qquad f : \prod_{s : S} \bigl(Q(u\,s) \to P s\bigr):
\]
forward on shapes, \emph{backward} on positions. Its extension sends
$(s, \varphi) \mapsto (u\,s,\; \varphi \circ f_s)$.
\end{definition}

\begin{theorem}[Representation theorem, AAG]\label{thm:aag}
The extension functor $\Cont \to [\Set,\Set]$ is full and faithful: container
morphisms are exactly the natural transformations between the extensions.
\end{theorem}

Containers are closed under $+$, $\times$, composition, $\Sigma$, $\Pi$ and
fixpoints, so every strictly positive type is a container; the derivative
$\partial(S \con P) = \bigl(\sum_s P s\bigr) \con \lambda(s,p).\,P s \setminus
\{p\}$ gives one-hole contexts and Huet's zippers.

\begin{remark}[What container morphisms cannot express]
A morphism supplies \emph{one} reindexing per shape, applied uniformly to every
labelling of that shape. Sorting chooses its permutation by inspecting values, so
it is not natural in $X$; by Theorem~\ref{thm:aag} no sorting algorithm is a
container morphism. This is a proof, not an intuition, and it is worth keeping in
mind as the guide to what the formalism does and does not see.
\end{remark}

\subsection{Containers are \texorpdfstring{$\FF_1$}{F1}-linear algebra}

Quotient containers, in which a group acts on the positions of each shape, are
Joyal's analytic functors --- species. The symmetry group of $n$ positions is
$S_n$. And $S_n = \GL_n(\FF_1)$ (FFD \S16). So:

\begin{center}
\fbox{\parbox{0.9\textwidth}{A shape with a \emph{set} of positions is a shape
with an $\FF_1$-vector space of positions. The positions have no addition, only
relabelling. Container theory is the type theory of the field with one element.}}
\end{center}

The question of this document is what happens at a general $q$.

\section{The obstruction: one cannot internalise na\"ively}\label{sec:obstruction}

\begin{warning}\label{warn:collapse}
There is no useful notion of ``container internal to $\Vect_{\FF_q}$''.
\end{warning}

By Gambino--Kock, containers are the polynomial functors of a locally cartesian
closed category, built from $\Sigma$, $\Pi$ and pullback. $\Vect$ is not locally
cartesian closed; it is additive, so finite products coincide with coproducts. In
$\Set$, $X^{P s}$ is a product of $\lvert P s\rvert$ copies of $X$; in $\Vect$
the product of $n$ copies of $V$ is $V^{\oplus n}$, so
\[
  \sum_s V^{\prod P s} \;=\; \bigoplus_s V^{\oplus \lvert P s\rvert}
\]
is \emph{linear} in $V$. Every polynomial degree above one is destroyed.

\begin{remark}
This is the topos/abelian dichotomy of FFD \S11: no nontrivial category is both a
topos and abelian. Containers are topos-side technology --- disjoint coproducts,
$\Sigma \dashv \Pi$, subobject lattices that are Boolean. $\Vect$ is on the
abelian side. The deformation must therefore change something other than the
ambient category, and identifying \emph{what} is the whole design problem.
\end{remark}

\section{The right formulation, and the definition}

Read the extension again and notice what it actually is. Writing
$\yon^{W} = \Hom(W,-)$ for the covariant functor corepresented by $W$,
\[
  \ext{S \con P} \;=\; \sum_{s : S} \yon^{P s}.
\]

\begin{center}
\fbox{\parbox{0.9\textwidth}{\textbf{A container is a set-indexed coproduct of
corepresentable functors.} Theorem~\ref{thm:aag} is then the Yoneda lemma: natural
transformations $\yon^{W} \Rightarrow \yon^{W'}$ are maps $W' \to W$, which is
exactly why container morphisms run backwards on positions.}}
\end{center}

That formulation says nothing about $\Set$ being the base. Change the base.

\begin{definition}[$q$-container]\label{def:qcont}
Fix a finite field $\FF_q$. A \emph{$q$-container} is a pair $S \con P$ with $S$ a
finite set of \emph{shapes} and $P : S \to \finVect_{\FF_q}$ assigning to each
shape a finite-dimensional \emph{position space}. Its \emph{extension} is the
functor $\finVect_{\FF_q} \to \Set$
\[
  \ext{S \con P} V \;=\; \sum_{s \in S} \Hom_{\FF_q}(P s,\, V).
\]
\end{definition}

\begin{definition}[$q$-container morphism]\label{def:qcontmor}
A morphism $(S \con P) \to (T \con Q)$ is a pair
\[
  u : S \to T, \qquad
  f \in \prod_{s \in S} \Hom_{\FF_q}\bigl(Q(u\,s),\, P s\bigr),
\]
the $f_s$ being \emph{linear} maps, backwards on positions. Composition is
composition of the $u$'s and (contravariantly) of the $f$'s. Write $\Cont_q$ for
the resulting category.
\end{definition}

Labelling a shape now means giving a \emph{linear map} from its position space
into $V$: positions can be superposed, scaled and linearly recombined, and
reindexing must respect that.

\begin{remark}[Skeleton]
$\finVect_{\FF_q}$ is skeletal up to isomorphism on dimensions, so a $q$-container
is determined up to isomorphism by the multiset of dimensions
$(\dim P s)_{s \in S}$. All computations below use this presentation.
\end{remark}

\begin{example}\label{ex:basic}
\begin{enumerate}[nosep,label=(\alph*)]
  \item $S = \{\ast\}$, $P\ast = \FF_q^{0}$: the constant functor $1$.
  \item $S = \{\ast\}$, $P\ast = \FF_q^{1}$: the identity, since
        $\Hom(\FF_q, V) \cong V$.
  \item $S = \{\ast\}$, $P\ast = \FF_q^{d}$: the functor $V \mapsto V^{\oplus d}$
        --- \emph{not} $V^{\otimes d}$ (see Warning~\ref{warn:collapse}).
  \item $S = \NN$, $P n = \FF_q^{n}$: the ``$q$-list'' container,
        $\ext{-}V = \coprod_n \Hom(\FF_q^n, V)$.
\end{enumerate}
\end{example}

\begin{proposition}[$q \to 1$]
Taking $\Vect_{\FF_1} = \FinSet_{\ast}$ with $\Hom_{\FF_1}$ the pointed maps
(FFD \S16), Definitions~\ref{def:qcont} and \ref{def:qcontmor} specialise to
pointed containers and pointed container morphisms.
\end{proposition}

\section{Theorem A: the representation theorem survives}

\begin{theorem}\label{thm:A}
The extension functor $\Cont_q \to [\finVect_{\FF_q}, \Set]$ is full and
faithful. Explicitly, with $d_s = \dim P s$ and $e_t = \dim Q t$,
\[
  \Cont_q(C,D) \;\cong\; \Nat\bigl(\ext{C},\ext{D}\bigr)
  \;\cong\; \prod_{s \in S} \; \sum_{t \in T} \Hom_{\FF_q}(Q t, P s),
\]
\[
  \text{so that}\qquad
  \bigl\lvert \Cont_q(C,D) \bigr\rvert \;=\; \prod_{s} \sum_{t} q^{\,e_t d_s}.
\]
\end{theorem}

\begin{proof}
Colimits in a functor category are pointwise, so $\Sigma$ is the coproduct in
$[\finVect_q, \Set]$ and
$\Nat\bigl(\sum_s F_s, G\bigr) \cong \prod_s \Nat(F_s, G)$. By Yoneda,
$\Nat(\yon^{W}, G) \cong G(W)$. Hence
\[
  \Nat(\ext{C},\ext{D})
  \;\cong\; \prod_{s} \ext{D}(P s)
  \;=\; \prod_{s} \sum_{t} \Hom(Q t, P s).
\]
On the other side, Definition~\ref{def:qcontmor} gives
$\Cont_q(C,D) = \sum_{u : S \to T} \prod_{s} \Hom(Q(u\,s), P s)$, and the
distributive law of $\Pi$ over $\Sigma$ in $\Set$,
\[
  \prod_{s \in S} \sum_{t \in T} A(s,t)
  \;\cong\; \sum_{u : S \to T} \; \prod_{s \in S} A\bigl(s, u\,s\bigr),
\]
identifies the two. Chasing the bijection shows it is exactly the extension map.
\end{proof}

\begin{remark}[Where the proof spends its budget]\label{rem:budget}
The proof uses precisely two facts about the \emph{indexing}: that $\Sigma$ is a
coproduct with disjoint, connected summands, and that $\Pi$ distributes over
$\Sigma$. Both are topos-side properties of $\Set$, and neither holds for $\oplus$
in $\Vect$. Theorem~\ref{thm:B} is the consequence.
\end{remark}

\subsection{Computational verification}

Verified by exact enumeration, two methods, neither assuming Yoneda:
\emph{exhaustive} search over all function families
$\{\alpha_n : \ext{C}(\FF_q^n) \to \ext{D}(\FF_q^n)\}$ keeping the natural ones;
and \emph{generic}, enumerating candidates by their value on
$(s, \mathrm{id}_{P s})$, verifying naturality of each in full, and separately
checking that every source element is $\ext{C}(\varphi)$ of a generic (so no
natural family can escape the count).

\begin{center}
\small
\setlength{\tabcolsep}{4.5pt}
\begin{tabular}{ccccrrc}
\toprule
$q$ & $C$ & $D$ & $\max\dim$ & method & $\#$natural & $\lvert\Cont_q(C,D)\rvert$ \\
\midrule
2 & $(1)$ & $(1)$ & 1 & exhaustive & 2 & 2 \\
2 & $(1)$ & $(1,1)$ & 1 & exhaustive & 4 & 4 \\
2 & $(1,1)$ & $(1)$ & 1 & exhaustive & 4 & 4 \\
2 & $(0,1)$ & $(1)$ & 1 & exhaustive & 2 & 2 \\
2 & $(1,1)$ & $(1,1)$ & 1 & exhaustive & 16 & 16 \\
3 & $(1)$ & $(1)$ & 1 & exhaustive & 3 & 3 \\
3 & $(1)$ & $(1,1)$ & 1 & exhaustive & 6 & 6 \\
2 & $(1)$ & $(2)$ & 2 & exhaustive & 4 & 4 \\
2 & $(2)$ & $(1)$ & 2 & generic & 4 & 4 \\
2 & $(2)$ & $(1,1)$ & 2 & generic & 8 & 8 \\
2 & $(1,2)$ & $(2,1)$ & 2 & generic & 120 & 120 \\
3 & $(2)$ & $(1)$ & 2 & generic & 9 & 9 \\
\bottomrule
\end{tabular}
\end{center}

\noindent
A separate elementwise check ($q=2$, $C = D = (1,1)$) confirms that the map from
container morphisms to natural families is injective and surjective, not merely
equinumerous.

\section{Theorem B: linearising the shapes destroys the theory}

Definition~\ref{def:qcont} keeps the extension $\Set$-valued. Since
$\Hom_{\FF_q}(P s, V)$ is itself a vector space, one is tempted to replace
$\Sigma$ by $\oplus$ and land in $\Vect$. This is a mistake, and the failure is
total rather than marginal.

\begin{definition}
The \emph{linear extension} of $S \con P$ is the $\FF_q$-linear functor
$\ext{S \con P}_{\mathrm{lin}} V = \bigoplus_{s \in S} \Hom_{\FF_q}(P s, V)$.
\end{definition}

\begin{theorem}\label{thm:B}
For finite $S$ and $T$,
\[
  \Nat\bigl(\ext{C}_{\mathrm{lin}}, \ext{D}_{\mathrm{lin}}\bigr)
  \;\cong\; \Hom_{\FF_q}\Bigl( \bigoplus_{t \in T} Q t,\; \bigoplus_{s \in S} P s \Bigr),
  \qquad
  \text{of cardinality } q^{\left(\sum_s d_s\right)\left(\sum_t e_t\right)}.
\]
In particular the right-hand side depends only on the \emph{total position
spaces}: the shape sets $S$ and $T$ are forgotten entirely.
\end{theorem}

\begin{proof}
A natural transformation between additive functors is automatically
$\FF_q$-linear in each component (use the scalar endomorphisms $\lambda\cdot\mathrm{id}$
for homogeneity, and the diagonal and codiagonal $V \to V \oplus V \to V$ for
additivity). Enriched Yoneda gives $\Nat(\Hom(W,-), G) \cong G(W)$ for linear $G$,
and $\oplus$ is a biproduct, hence also a product, so
\[
  \Nat\Bigl(\bigoplus_s \Hom(P s,-), G\Bigr)
  \;\cong\; \prod_s G(P s)
  \;=\; \prod_s \bigoplus_t \Hom(Q t, P s)
\]
\[
  \phantom{\Nat}\;\cong\; \prod_s \Hom\Bigl(\bigoplus_t Q t,\, P s\Bigr)
  \;\cong\; \Hom\Bigl(\bigoplus_t Q t,\; \bigoplus_s P s\Bigr). \qedhere
\]
\end{proof}

\begin{corollary}
The linear extension functor is \textbf{neither full nor faithful}.
\end{corollary}

\begin{example}[Both failures at once, and cardinality is no guide]
Take $q = 2$, $C = (1)$, $D = (1,1)$. Then
$\lvert\Cont_q(C,D)\rvert = 4$ and $\lvert\Nat_{\mathrm{lin}}\rvert = 4$ --- the
cardinalities agree --- yet the map between them is neither injective nor
surjective, having image of size $3$:
\begin{itemize}[nosep]
  \item \emph{not faithful}: the two morphisms $(u = 0, f = 0)$ and
        $(u = 1, f = 0)$ have the same linear extension, namely $0$. In $\Set$
        these land in different summands and stay apart; in $\Vect$ the zero
        vector is shared.
  \item \emph{not full}: the ``diagonal'' map $(1,1) \in \Hom(\FF_2^{2}, \FF_2)$
        is natural but is not the extension of any container morphism, since a
        morphism must choose a single target shape per source shape.
\end{itemize}
\end{example}

\begin{center}
\small
\begin{tabular}{ccrrrcc}
\toprule
$q$ & $C \to D$ & $\lvert\Cont_q\rvert$ & image & $\lvert\Nat_{\mathrm{lin}}\rvert$ & faithful & full \\
\midrule
2 & $(1) \to (1)$ & 2 & 2 & 2 & yes & yes \\
2 & $(1) \to (1,1)$ & 4 & 3 & 4 & \NO & \NO \\
2 & $(1,1) \to (1,1)$ & 16 & 9 & 16 & \NO & \NO \\
2 & $(1) \to (1,1,1)$ & 6 & 4 & 8 & \NO & \NO \\
3 & $(1) \to (1,1)$ & 6 & 5 & 9 & \NO & \NO \\
\bottomrule
\end{tabular}
\end{center}

\begin{center}
\fbox{\parbox{0.92\textwidth}{\textbf{The representation theorem survives exactly
as long as the shapes stay a set.} Positions may be linearised freely --- that is
Theorem~\ref{thm:A}. Linearising the shapes replaces a disjoint union by a
biproduct, natural transformations may then take linear combinations across
different target shapes, and by Theorem~\ref{thm:B} the shape set is annihilated.
This is Remark~\ref{rem:budget} made quantitative, and the topos/abelian
dichotomy of FFD \S11 in the form of a cardinality identity.}}
\end{center}

\section{Counting, and the Jackson \texorpdfstring{$q$}{q}-derivative}

\begin{proposition}[Counting series]\label{prop:count}
$\bigl\lvert \ext{S \con P}(\FF_q^{n}) \bigr\rvert = \sum_{s} \bigl(q^{n}\bigr)^{d_s}$.
Writing $x = q^{n}$, the counting series of a $q$-container is
$f_C(x) = \sum_{s} x^{d_s}$, exactly parallel to
$\lvert\ext{S \con P}(X)\rvert = \sum_s \lvert X\rvert^{\lvert P s\rvert}$ in
$\Set$, under $\lvert X \rvert \leftrightarrow q^{n}$ and
$\lvert P s\rvert \leftrightarrow \dim P s$.
\end{proposition}

Verified by direct enumeration for $q = 2$, $n \le 3$, $C \in \{(1),(2),(1,2)\}$.

\begin{definition}[Derivative]\label{def:deriv}
For a $q$-container $S \con P$, put
\[
  \partial(S \con P) \;=\; \Bigl(\textstyle\sum_{s} \{\text{hyperplanes of } P s\}\Bigr)
  \;\con\; \bigl((s,H) \mapsto H\bigr),
\]
so a shape of $\partial C$ is a shape of $C$ together with a hyperplane of its
position space, and the labelling forgets the missing direction.
\end{definition}

\begin{definition}[Jackson $q$-derivative]
$D_q f(x) = \dfrac{f(qx) - f(x)}{(q-1)x}$, so that
$D_q x^{d} = [d]_q\, x^{d-1}$ with $[d]_q = 1 + q + \cdots + q^{d-1}$.
\end{definition}

\begin{theorem}\label{thm:D}
$f_{\partial C} = D_q f_C$. As $q \to 1$, $[d]_q \to d$, so $D_q \to d/dx$ and
$\partial$ degenerates to the ordinary container derivative.
\end{theorem}

\begin{proof}
$\FF_q^{d}$ has $\qbin{d}{d-1} = [d]_q$ hyperplanes, each of dimension $d-1$.
\end{proof}

Verified exactly for $q \in \{2,3\}$ and
$C \in \{(1),(2),(3),(1,2),(2,2,3)\}$; e.g.\ $q = 3$, $C = (2,2,3)$ gives
$f_C = x^{3} + 2x^{2}$ and $f_{\partial C} = 13x^{2} + 8x = D_q f_C$.

\section{A refuted conjecture, and what it was really about}\label{sec:refuted}

\begin{refuted}\label{conj:two}
We conjectured that a $q$-container has \emph{two} derivatives ---
$\partial_{\mathrm{sub}}$, restricting to a hyperplane of $P s$, and
$\partial_{\mathrm{quot}}$, quotienting by a line --- exchanged by the polarity
$W \mapsto W^{\perp}$ of FFD \S14, and merging at $q = 1$. The evidence offered
was that the Gaussian Pascal recursion has two forms,
\[
  \qbin{n}{k} = \qbin{n-1}{k-1} + q^{k}\qbin{n-1}{k},
  \qquad
  \qbin{n}{k} = q^{\,n-k}\qbin{n-1}{k-1} + \qbin{n-1}{k},
\]
which coincide at $q = 1$.
\end{refuted}

\begin{proposition}[Refutation]
Over $\FF_q$, $\partial_{\mathrm{sub}} \cong \partial_{\mathrm{quot}}$.
\end{proposition}

\begin{proof}
$\qbin{d}{d-1} = \qbin{d}{1} = [d]_q$, so a shape has as many hyperplanes as
lines; and $H \cong P s / \ell \cong \FF_q^{\,d-1}$ for every hyperplane $H$ and
line $\ell$. Since a $q$-container remembers only the dimension of each position
space, the two constructions give isomorphic $q$-containers.
\end{proof}

Confirmed computationally for $q \in \{2,3\}$ and all position dimensions
$\le 3$. So the two Pascal recursions are not two derivatives; they are
\emph{one derivative counted two ways}.

\begin{remark}[The diagnosis]
The conjecture silently assumed a base category in which subobjects and quotients
are not interchangeable. $\Vect_{\FF_q}$ is semisimple: every short exact
sequence splits, and the category is self-dual. The interesting deformation is
therefore not $q$ but the \emph{failure of self-duality of the base}. Note that
this is a different failure of splitting from the one in FFD Part~I: there, a
\emph{form} on a semisimple category fails to split $V = W \perp W^{\perp}$; here
it is the \emph{category itself} that splits too well. Conflating them was the
error.
\end{remark}

\begin{conjecture}[Repaired]\label{conj:repaired}
Let $\mathcal{A}$ be a finite $\FF_q$-linear abelian category and define
$\mathcal{A}$-containers by Definition~\ref{def:qcont} with
$P : S \to \mathcal{A}$. Then $\partial_{\mathrm{sub}} \cong
\partial_{\mathrm{quot}}$ if and only if $\mathcal{A}$ is self-dual. In
particular they separate over $\mathcal{A} = \Rep_{\FF_q}(\vec{A_n})$.
\end{conjecture}

\begin{example}[Smallest witness, hand-checked]
Let $\mathcal{A} = \Rep_{\FF_q}(1 \to 2)$, with simples $S_1 = (\FF_q, 0)$,
$S_2 = (0, \FF_q)$ and projective $P_1 = (\FF_q \xrightarrow{\ \mathrm{id}\ } \FF_q)$.
The proper nonzero subobjects of $P_1$ are just $S_2$, since a subrepresentation
$(W_1, W_2)$ needs $\mathrm{id}(W_1) \subseteq W_2$; and $P_1/S_2 = S_1$. As
$S_1 \not\cong S_2$,
\[
  \partial_{\mathrm{sub}}(P_1) = S_2 \;\ne\; S_1 = \partial_{\mathrm{quot}}(P_1).
\]
\end{example}

\begin{remark}
One might expect the repair to the derivative and the enrichment of the
composition (\S\ref{sec:hall}) to have a single cause, namely non-semisimplicity.
They do not. Non-semisimplicity enriches the Hall algebra, but a category can be
non-semisimple and still self-dual --- finite modules over a discrete valuation
ring are, by Matlis duality --- in which case the two derivatives still coincide.
The two phenomena sit on independent axes, tabulated in
Table~\ref{tab:threebases}.
\end{remark}

\section{Composition, and where Hall--Littlewood enters}\label{sec:hall}

\emph{The classical theorems quoted here are Hall's and Macdonald's; their
application to $q$-containers is design, and nothing in this section has been
verified computationally. An earlier draft of this section asserted that Day
convolution has Hall-algebra structure constants. That is false, and
\S\ref{sec:diverge} is the correction.}

\subsection{Two candidate composition laws}

\begin{definition}[Day convolution]\label{def:day}
Let $(\mathcal{A}, \otimes)$ be a small monoidal category. The functor category
$[\mathcal{A}, \Set]$ carries the monoidal structure
\[
  (F \star G)(c) \;=\; \int^{a,b} \mathcal{A}(a \otimes b,\, c) \times F(a) \times G(b),
\]
the unique one making the Yoneda embedding strong monoidal.
\end{definition}

\begin{example}[The species product]
For $\mathcal{A} = (\FinSet_{\mathrm{bij}}, \sqcup)$ the coend collapses to
$(F \star G)[n] = \sum_{A \sqcup B = [n]} F[A] \times G[B]$, the species product,
which decategorifies to multiplication of exponential generating functions.
\end{example}

\begin{definition}[Hall algebra]\label{def:hall}
Let $\mathcal{A}$ be an abelian category in which every object has finitely many
subobjects and all $\Hom$ and $\mathrm{Ext}^1$ sets are finite. Its \emph{Hall
algebra} $H(\mathcal{A})$ is the free $\ZZ$-module on isomorphism classes, with
\[
  [M]\cdot[N] \;=\; \sum_{[L]} g^{L}_{MN}\,[L],
  \qquad
  g^{L}_{MN} \;=\; \#\bigl\{\, U \subseteq L \;:\; U \cong N,\; L/U \cong M \,\bigr\}.
\]
So $g^{L}_{MN}$ counts \emph{subobjects with prescribed quotient} --- that is, all
extensions of $M$ by $N$ with middle term $L$, split or not.
\end{definition}

These are genuinely different operations, and it is worth computing both in the
smallest case before relying on either.

\begin{proposition}[Day convolution over $\FF_q$]\label{prop:day}
Let $\mathcal{G}_q = \bigsqcup_n B\GL_n(\FF_q)$ be the groupoid of
finite-dimensional $\FF_q$-spaces and isomorphisms, monoidal under $\oplus$, and
let $\delta_k$ be the point-functor supported on dimension $k$. Then
\[
  (\delta_k \star \delta_{n-k})(\FF_q^{n})
  \;=\; \GL_n \big/ \bigl(\GL_k \times \GL_{n-k}\bigr),
  \qquad
  \text{of cardinality } \qbin{n}{k}\, q^{\,k(n-k)} :
\]
the number of \emph{ordered direct-sum decompositions} of $\FF_q^{n}$ into
subspaces of dimensions $k$ and $n-k$.
\end{proposition}

\begin{example}
$n = 2$, $k = 1$, $q = 2$: the formula gives $3 \cdot 2 = 6$, and directly,
$\FF_2^{2}$ has $3$ lines, each with $2$ complements.
\end{example}

\begin{proposition}[The Hall product over $\FF_q$]\label{prop:thin}
For $\mathcal{A} = \Vect_{\FF_q}$, $g^{\,n}_{k,\,n-k} = \qbin{n}{k}$, the number of
subspaces of the given dimension. In the same case as above this is $3$, not $6$.
\end{proposition}

\begin{corollary}[Agreement, up to a twist, when the base is semisimple]
Over $\Vect_{\FF_q}$ the two structure constants differ by the factor
$q^{\,k(n-k)} = q^{\langle M, N\rangle}$, where
$\langle M,N\rangle = \dim\Hom(M,N) - \dim\mathrm{Ext}^1(M,N)$ is the Euler form
--- the number of complements to a fixed subspace. This is exactly the twist
relating the Hall algebra to its twisted form, so over a semisimple base the
discrepancy is cosmetic.
\end{corollary}

\subsection{Over a filtered base they diverge essentially}\label{sec:diverge}

\begin{definition}[Discrete valuation ring, type of a module]
A \emph{discrete valuation ring} $\mathcal{O}$ is a local principal ideal domain
that is not a field; write $\mathfrak{m} = (\pi)$ and assume the residue field
$\mathcal{O}/\mathfrak{m} \cong \FF_q$ is finite. Canonical examples:
$\FF_q[[t]]$, and $\ZZ_p$ with $q = p$. Every finite $\mathcal{O}$-module is
uniquely $M_\lambda \cong \bigoplus_i \mathcal{O}/(\pi^{\lambda_i})$ for a
partition $\lambda$, its \emph{type}; write $n(\lambda) = \sum_i (i-1)\lambda_i$.
\end{definition}

Such a module is an $\FF_q$-space with a nilpotent operator: positions carrying a
filtration, a ``shift with decay''. Setting $\pi = 0$ recovers
$\Vect_{\FF_q}$, so this enriches Definition~\ref{def:qcont} rather than
departing from it.

\begin{example}[The computation that separates them]\label{ex:separate}
Multiply the class of the simple module by itself.
\[
  \textbf{Hall:}\qquad
  u_{(1)} \cdot u_{(1)} \;=\; u_{(2)} \;+\; (q+1)\, u_{(1,1)} .
\]
The term $u_{(2)}$ arises from the \emph{non-split} extension
$\mathcal{O}/\pi \subset \mathcal{O}/\pi^{2}$; the coefficient $q+1$ counts the
lines in $\FF_q^{2}$.
\[
  \textbf{Day convolution along } \oplus:\qquad
  \delta_{(1)} \star \delta_{(1)} \;=\; q(q+1)\, \delta_{(1,1)},
  \quad\text{and nothing else.}
\]
\end{example}

\begin{proposition}\label{prop:daymisses}
$(F \star G)(L)$ is empty unless $L \cong U \oplus W$ for some $U, W$ in the
supports of $F$ and $G$. Hence Day convolution along $\oplus$ cannot see
non-split extensions at all.
\end{proposition}

\begin{center}
\fbox{\parbox{0.92\textwidth}{\textbf{Correction.} Day convolution along $\oplus$
is the Cauchy product of $q$-species; its structure constants count
\emph{splittings}, and agree with the twisted Hall constants only when the base is
semisimple. It is \emph{not} the Hall product. Reaching the Hall algebra requires
convolution along the span of short exact sequences,
\[
  \mathcal{A} \times \mathcal{A} \;\longleftarrow\; \mathcal{E} \;\longrightarrow\; \mathcal{A},
  \qquad
  \bigl(0 \to N \to L \to M \to 0\bigr) \;\longmapsto\; \bigl((M,N),\, L\bigr),
\]
pushing forward along the middle-term map --- the Ringel--Hall product as a
convolution algebra on the stack of objects (Kapranov--Vasserot; Schiffmann's
lectures). Day convolution deletes precisely the term that carries the content.

\medskip
\noindent\S\ref{sec:indres} identifies what the two products actually are:
convolution along $\oplus$ induces from the \emph{Levi}, the Hall product induces
from the \emph{parabolic}, and the discrepancy is the unipotent radical.}}
\end{center}

\subsection{The Hall side is Hall--Littlewood}

\begin{definition}[Hall--Littlewood polynomials]
$P_\lambda(x;t)$ is the basis of the ring $\Lambda$ of symmetric functions that is
monic--triangular against the monomial basis for dominance order and orthogonal
for the $t$-deformed Hall inner product. It interpolates: $P_\lambda(x;0) =
s_\lambda$ and $P_\lambda(x;1) = m_\lambda$.
\end{definition}

\begin{theorem}[Hall; Macdonald, \emph{Symmetric Functions and Hall Polynomials},
II.4 and III.3]\label{thm:HL}
Let $\mathcal{A}_{\mathcal{O}}$ be the category of finite modules over a discrete
valuation ring with residue field $\FF_q$. The structure constants
$g^{\lambda}_{\mu\nu}(q)$ of Definition~\ref{def:hall} are polynomials in $q$ with
integer coefficients --- the \emph{Hall polynomials} --- independent of which
$\mathcal{O}$ realises the residue field, and
\[
  H(\mathcal{A}_{\mathcal{O}}) \;\xrightarrow{\ \sim\ }\; \Lambda,
  \qquad
  u_\lambda \;\longmapsto\; q^{-n(\lambda)}\,P_\lambda\!\left(x;\,q^{-1}\right)
\]
is a ring isomorphism. Equivalently
$f^{\lambda}_{\mu\nu}(t) = t^{\,n(\lambda)-n(\mu)-n(\nu)} g^{\lambda}_{\mu\nu}(t^{-1})$
are the Hall--Littlewood structure constants.
\end{theorem}

\begin{example}[Example~\ref{ex:separate} pushed through]
With $n((1)) = n((2)) = 0$ and $n((1,1)) = 1$, the Hall identity becomes
\[
  P_1^{\,2} \;=\; P_{(2)} \;+\; \bigl(1 + q^{-1}\bigr) P_{(1,1)}
  \;=\; P_2 + (1+t)\,P_{11}.
\]
At $t = 0$ this is $s_1^{2} = s_2 + s_{11}$; at $t = 1$ it is
$m_1^{2} = m_2 + 2m_{11}$. Both correct. Note that the content sits in the
$P_2$ term --- the non-split extension --- which is exactly what
Proposition~\ref{prop:daymisses} deletes.
\end{example}

\begin{center}
\fbox{\parbox{0.92\textwidth}{\textbf{Weakened conclusion.} $q$-containers with
$\mathcal{O}$-module positions admit \emph{two} candidate composition laws. Day
convolution along $\oplus$ gives the Cauchy product; convolution along short exact
sequences gives the Hall product and hence, by Theorem~\ref{thm:HL},
multiplication of Hall--Littlewood polynomials $P_\lambda(x;q^{-1})$. The two
agree up to the Euler twist over $\Vect_{\FF_q}$ and disagree essentially the
moment the positions carry a filtration. Which one deserves the name
``composition of containers'' is open --- see
Question~\ref{q:closed}.}}
\end{center}

\begin{remark}[Green's polynomials close the circle]
The same $P_\lambda(x;q^{-1})$ carry the character theory of the groups whose
$\FF_1$-shadow started this document:
$p_\mu = \sum_\lambda Q^{\lambda}_{\mu}(q)\,P_\lambda(x;q^{-1})$ with
$Q^{\lambda}_{\mu}(q) \in \ZZ[q]$ the unipotent character values of
$\GL_n(\FF_q)$, and $Q^{\lambda}_{\mu}(1) = \chi^{\lambda}_{\mu}$, the character
table of $S_n$. Composing with Kostka--Foulkes,
$s_\lambda = \sum_\nu K_{\lambda\nu}(q)P_\nu(x;q)$, factors this as
$Q = \chi^{\mathsf{T}}K$. So the Hall--Littlewood basis is the coordinate system
in which Zelevinsky's tower $\bigoplus_n \Rep(\GL_n(\FF_q))$ --- degenerating to
$\bigoplus_n \Rep(S_n)$, hence to species --- is legible.
\end{remark}

\subsection{Two independent axes}

It is tempting to say that ``non-semisimple base'' is the single repair that both
enriches the composition and separates the two derivatives of
\S\ref{sec:refuted}. That is false, and the distinction matters.

\begin{definition}[Self-dual category]
$\mathcal{A}$ is \emph{self-dual} if there is an exact anti-equivalence
$\mathcal{A}^{\op} \simeq \mathcal{A}$, that is, a duality exchanging subobjects
and quotients. $\Vect_{\FF_q}$ is self-dual by $V \mapsto V^{*}$;
$\mathcal{A}_{\mathcal{O}}$ is self-dual by Matlis duality
$M \mapsto \Hom_{\mathcal{O}}(M, K/\mathcal{O})$, which preserves the type
$\lambda$; $\Rep_{\FF_q}(\vec{A_2})$ is \emph{not}, as $S_1 \not\cong S_2$
witnesses.
\end{definition}

\begin{table}[h]
\centering
\small
\setlength{\tabcolsep}{5pt}
\begin{tabular}{lccc}
\toprule
base $\mathcal{A}$ for positions & Hall algebra & self-dual? & $\partial_{\mathrm{sub}}$ vs $\partial_{\mathrm{quot}}$ \\
\midrule
$\Vect_{\FF_q}$
  & divided powers, $\qbin{n}{k}$ & yes & coincide (\S\ref{sec:refuted}, proved) \\
finite $\mathcal{O}$-modules
  & $\Lambda$, Hall--Littlewood $P_\lambda(x;q^{-1})$ & yes (Matlis) & coincide (predicted) \\
$\Rep_{\FF_q}(\vec{A_2})$
  & $U_q(\mathfrak{n}^{-})$ for $\mathfrak{sl}_2$, noncomm. & \textbf{no} & \textbf{separate} (Conj.~\ref{conj:repaired}) \\
\bottomrule
\end{tabular}
\caption{Filtration depth and failure of self-duality are independent. The middle
row is non-semisimple, has a rich commutative Hall algebra, and still has
coincident derivatives; only the bottom row separates them.}
\label{tab:threebases}
\end{table}

\noindent
So: \emph{filtration depth} buys Hall--Littlewood; \emph{failure of self-duality}
buys two derivatives. Neither implies the other.

\subsection{What is now open, and what to compute}

\begin{question}\label{q:closed}
$q$-containers are the coproducts of corepresentables. Is that class closed under
Hall convolution? Day convolution along $\oplus$ preserves it, because
$\yon^{U} \star \yon^{W} \cong \yon^{U \oplus W}$. The Hall product has no such
easy answer, and if the class is not closed then ``Hall--Littlewood composition''
is a statement about $q$-species, not about $q$-containers.
\end{question}

\begin{enumerate}[nosep]
  \item Settle Question~\ref{q:closed}. This is the load-bearing one: it decides
        whether \S\ref{sec:hall} is about containers at all.
  \item Enumerate both products for $\mathcal{O} = \FF_2[t]/(t^{k})$, $k \le 3$,
        and confirm Example~\ref{ex:separate} and the Hall polynomials
        $g^{\lambda}_{\mu\nu}(q)$ for $\lvert\lambda\rvert \le 4$.
  \item Confirm the middle row of Table~\ref{tab:threebases}: that Matlis duality
        alone forces $\partial_{\mathrm{sub}} \cong \partial_{\mathrm{quot}}$,
        isolating failure of self-duality as the sole obstruction.
  \item Check that Theorem~\ref{thm:A} is stable under change of base; the proof
        used only that the shapes form a set.
\end{enumerate}


\section{Induction, restriction, and change of base}\label{sec:indres}

The two most useful operations in representation theory are induction and
restriction, and they are adjoint. This section shows that any such adjunction is
already native to the container calculus, and that doing so names the discrepancy
of \S\ref{sec:diverge}.

\subsection{Adjunctions lift, with a variance flip}

\begin{definition}[Change of base]
For a functor $A : \mathcal{A} \to \mathcal{B}$ and a container $C = S \con P$ with
positions in $\mathcal{A}$, write $A_{*}C = S \con (A \circ P)$, a container with
positions in $\mathcal{B}$. On morphisms, $A_{*}(u, f) = (u, A f)$.
\end{definition}

\begin{theorem}\label{thm:lift}
Let $A : \mathcal{A} \to \mathcal{B}$ and $B : \mathcal{B} \to \mathcal{A}$ with
$B \dashv A$. Then $A_{*} \dashv B_{*}$ on containers:
\[
  \Cont_{\mathcal{B}}\bigl(A_{*}C,\, D\bigr)
  \;\cong\;
  \Cont_{\mathcal{A}}\bigl(C,\, B_{*}D\bigr).
\]
\end{theorem}

\begin{proof}
By Theorem~\ref{thm:A} both sides are products of sums of hom-sets:
\begin{align*}
  \Cont_{\mathcal{B}}(A_{*}C, D) &= \prod_{s}\sum_{t} \Hom_{\mathcal{B}}\bigl(Q t,\, A(P s)\bigr), \\
  \Cont_{\mathcal{A}}(C, B_{*}D) &= \prod_{s}\sum_{t} \Hom_{\mathcal{A}}\bigl(B(Q t),\, P s\bigr),
\end{align*}
and $\Hom_{\mathcal{B}}(X, A Y) \cong \Hom_{\mathcal{A}}(B X, Y)$ is exactly
$B \dashv A$. The bijection is natural in each factor, hence in $C$ and $D$.
\end{proof}

\begin{remark}[Why the variance flips]
The \emph{left} adjoint on positions becomes the \emph{right} adjoint on
containers, and conversely. This is forced by Definition~\ref{def:qcontmor}:
container morphisms run backwards on positions, so an adjunction used on the
position category is consumed in the opposite slot.
\end{remark}

\begin{proposition}[Extensions]\label{prop:extadj}
If $A \dashv B$ then $\ext{A_{*}C} = \ext{C} \circ B$. Pushing positions along a
left adjoint is precomposing the extension with its right adjoint.
\end{proposition}

\begin{proof}
$\ext{A_{*}C}(W) = \sum_s \Hom_{\mathcal{B}}(A(P s), W) \cong \sum_s
\Hom_{\mathcal{A}}(P s, B W) = \ext{C}(B W)$.
\end{proof}

\subsection{Induction and restriction}

\begin{definition}
For $H \leq G$ and a field $k$: $\Res^G_H : \Rep_k(G) \to \Rep_k(H)$ restricts the
action; $\Ind_H^G = k[G] \otimes_{k[H]} - $ and $\Coind_H^G = \Hom_{k[H]}(k[G], -)$.
One has $\Ind \dashv \Res \dashv \Coind$, and for $[G:H] < \infty$,
$\Ind \cong \Coind$, so $\Res$ is \emph{biadjoint} to $\Ind$.
\end{definition}

\begin{corollary}\label{cor:indres}
For $H \leq G$ of finite index, positions in $\Rep_k(-)$ give
\[
  \Ind_{*} \;\dashv\; \Res_{*} \;\dashv\; \Ind_{*},
\]
and, by Proposition~\ref{prop:extadj},
\[
  \ext{\Ind_{*}C} = \ext{C} \circ \Res,
  \qquad
  \ext{\Res_{*}C} = \ext{C} \circ \Ind .
\]
\end{corollary}

Inducting the positions is restricting the values, and conversely: Frobenius
reciprocity holds shape by shape. Counting, with $k = \FF_q$,
\[
  \bigl\lvert\Cont_{G}(C,D)\bigr\rvert
  \;=\; \prod_{s}\sum_{t} q^{\,\dim \Hom_{G}(Q t,\, P s)},
\]
and in the semisimple case $\dim\Hom_G(M,N) = \langle \chi_M, \chi_N\rangle$: the
container hom-set is computed by inner products of characters. Taking $G$ trivial
recovers Theorem~\ref{thm:A}.

\subsection{Parabolic versus Levi: what \S\ref{sec:diverge} was about}

Let $G = \GL_n(\FF_q)$, let $P$ be the parabolic stabilising a $k$-dimensional
subspace, $L = \GL_k \times \GL_{n-k}$ its Levi factor and $U$ its unipotent
radical, so $P = L \ltimes U$.

\begin{definition}[Harish-Chandra induction]
$R_L^G = \Ind_P^G \circ \Infl_L^P : \Rep(L) \to \Rep(G)$: inflate along
$P \twoheadrightarrow L$, killing $U$, then induce. Its adjoint $^{*}R^G_L$ takes
$U$-invariants and restricts to $L$.
\end{definition}

\begin{proposition}\label{prop:parlevi}
\[
  [G:P] = \qbin{n}{k}, \qquad
  [G:L] = \qbin{n}{k} q^{\,k(n-k)}, \qquad
  [P:L] = \lvert U\rvert = q^{\,k(n-k)} .
\]
\end{proposition}

\begin{center}
\small
\setlength{\tabcolsep}{5pt}
\begin{tabular}{lccl}
\toprule
operation & induces from & index & container side \\
\midrule
Harish-Chandra $R_L^G = \Ind_P^G \Infl$ & parabolic $P$ & $\qbin{n}{k}$ & Hall product \\
ordinary $\Ind_L^G$ & Levi $L$ & $\qbin{n}{k}q^{k(n-k)}$ & Day conv.\ along $\oplus$ \\
\bottomrule
\end{tabular}
\end{center}

\begin{center}
\fbox{\parbox{0.92\textwidth}{\textbf{The discrepancy of \S\ref{sec:diverge} is the
unipotent radical.} Day convolution along $\oplus$ induces from the \emph{Levi}: it
remembers a choice of complement, which is not canonical, and the redundancy is
exactly $\lvert U\rvert = q^{k(n-k)}$. Harish-Chandra induction goes through the
\emph{parabolic} and inflates $U$ away, which is what lets it see extensions rather
than splittings. The factor I could previously only describe as an Euler-form twist
is the order of the unipotent radical.}}
\end{center}

\begin{example}[$\GL_2(\FF_2)$]
$\lvert G\rvert = 6$; $L$ is trivial so $[G:L] = 6$, the ordered pairs of
complementary lines; $P = B$ has order $2$, so $[G:P] = 3$, the lines; and
$\lvert U\rvert = 2 = q^{1}$.
\end{example}

\begin{remark}[Zelevinsky, and $q = 1$]
Taking induction and restriction to be \emph{Harish-Chandra} induction and
restriction along Levi factors and their unipotent subgroups is exactly what gives
$\bigoplus_n \Rep(\GL_n(\FF_q))$ its structure as a positive self-adjoint Hopf
algebra, decomposing as a tensor product of copies of the elementary one indexed by
cuspidal representations (Zelevinsky). At $q = 1$ there is no unipotent radical to
inflate away, the two rows of the table coincide, and one recovers the induction
product $\Ind_{S_k \times S_{n-k}}^{S_n}$ on $\bigoplus_n \Rep(S_n)$ --- which is
precisely the species product, i.e.\ Day convolution along $\sqcup$. The
parabolic/Levi distinction is a $q$-phenomenon, and it vanishes for the same reason
the two derivatives of \S\ref{sec:refuted} merge.
\end{remark}

\begin{remark}[The Hecke algebra]
$\Ind_B^G \mathbf{1} = k[G/B] = k[\text{complete flags}]$ is a free module on a
hom-set --- a container in the $k[-]$ sense of \S\ref{sec:hall} --- and
$\End_G(\Ind_B^G \mathbf{1}) = H_n(q)$, the Hecke algebra, degenerating to $k[S_n]$
at $q = 1$. So the Hecke algebra is the endomorphism algebra of the flag container.
\end{remark}

\begin{question}[Mackey]\label{q:mackey}
The Mackey formula $\Res \circ \Ind = \bigoplus_{HgK} \Ind \circ \operatorname{conj}
\circ \Res$ should lift to a distributive law between $\Ind_{*}$ and $\Res_{*}$. It
is the compatibility that makes the PSH structure work, so if it lifts cleanly the
container calculus carries both operations \emph{and} their interaction.
\end{question}


\section{A general programme: lifting Macdonald theory to containers}\label{sec:programme}

\subsection{Decategorification}

Proposition~\ref{prop:count} sends a $q$-container to a polynomial. That map is the
bottom of a ladder:
\begin{align*}
  &\underbrace{\text{$q$-containers}}_{\text{functors, natural bijections}}
  \;\longrightarrow\;
  \underbrace{\text{symmetric functions}}_{\text{characters, if the positions carry a group action}} \\[0.4em]
  &\qquad\longrightarrow\;
  \underbrace{\text{$q$-series}}_{f_C(x) = \sum_s x^{\dim P s}}
  \;\longrightarrow\;
  \text{numbers}
\end{align*}
and it is a homomorphism for every operation we have:

\begin{center}
\small
\setlength{\tabcolsep}{5pt}
\begin{tabular}{@{}>{\raggedright\arraybackslash}p{0.46\textwidth}>{\raggedright\arraybackslash}p{0.46\textwidth}@{}}
\toprule
container level & polynomial level \\
\midrule
coproduct $C + D$ & $f_C + f_D$ \\
product $C \times D$ (Def.~\ref{def:prod}) & $f_C \cdot f_D$ \\
Day convolution along $\oplus$ (Levi) & multiplication weighted by $q^{k(n-k)}$ \\
Hall product (parabolic, \S\ref{sec:indres}) & Gaussian-binomial structure constants \\
derivative $\partial$ (hyperplanes) & Jackson $D_q$ \\
shift $\sigma$ (Def.~\ref{def:shift}) & $f(x) \mapsto f(qx)$ \\
$\Ind_{*}$, $\Res_{*}$ (Cor.~\ref{cor:indres}) & multiplication, comultiplication \\
$q \to 1$ & $d/dx$, binomials, ordinary generating functions \\
\bottomrule
\end{tabular}
\end{center}

When the positions carry a group action the middle rung is not one-variable
polynomials but the ring $\Lambda$ of symmetric functions, and decategorification is
the character map: the Frobenius characteristic for $\bigoplus_n \Rep(S_n)$, and
Zelevinsky's for $\bigoplus_n \Rep(\GL_n(\FF_q))$, in which Green's polynomials give
the transition to the Hall--Littlewood basis (\S\ref{sec:hall}).

\subsection{The programme}

\begin{center}
\fbox{\parbox{0.92\textwidth}{\textbf{Programme.} Every identity in the $q$- and
$(q,t)$-theory of symmetric functions should be the shadow of a natural isomorphism
of containers, and every container construction should decategorify to an identity.
This is Joyal's species programme run at $q$: replace ``bijection of finite sets''
by ``isomorphism of $q$-containers'' and ``generating function'' by ``$q$-series''.
The payoff sought is not new identities but \emph{reasons} --- a coefficient that
looks arbitrary at the polynomial level should be forced at the container level.}}
\end{center}

\subsection{A worked instance: the \texorpdfstring{$q$}{q}-Leibniz rule}

\begin{definition}[Product of containers]\label{def:prod}
$(S \con P) \times (T \con Q) = (S \times T) \con \bigl((s,t) \mapsto P s \oplus Q t\bigr)$,
so that $\ext{C \times D} \cong \ext{C} \times \ext{D}$ pointwise and
$f_{C \times D} = f_C f_D$.
\end{definition}

\begin{definition}[Shift]\label{def:shift}
$\sigma C$ is the container with $\ext{\sigma C}V = \ext{C}(V \oplus \FF_q)$.
Explicitly $\sigma(S \con P) = \bigl(\sum_s (P s)^{*}\bigr) \con \bigl((s,\varphi)
\mapsto P s\bigr)$: each shape $s$ is repeated $q^{d_s}$ times, positions unchanged.
Hence $f_{\sigma C}(x) = f_C(qx)$, and $\sigma$ preserves containers.
\end{definition}

The Jackson derivative does not satisfy the ordinary Leibniz rule; it satisfies
\[
  D_q(fg) \;=\; D_q(f)\,g \;+\; \sigma(f)\,D_q(g),
  \qquad \sigma(f)(x) = f(qx),
\]
an asymmetry that is unmotivated at the level of polynomials. It is forced at the
level of containers.

\begin{theorem}[Categorified $q$-Leibniz]\label{thm:leibniz}
For $q$-containers $C, D$,
\[
  \partial(C \times D) \;\cong\; \partial C \times D \;+\; \sigma C \times \partial D .
\]
\end{theorem}

\begin{proof}
Both $\partial$ and $\times$ distribute over $+$, so it suffices to take
$C$ and $D$ with one shape each, of position dimensions $d$ and $e$. Then
$C \times D$ has the single shape of dimension $d+e$, so the left side has
$\qbin{d+e}{d+e-1} = [d+e]_q$ shapes, all of position dimension $d+e-1$. On the
right, $\partial C \times D$ contributes $[d]_q$ shapes and $\sigma C \times
\partial D$ contributes $q^{d}[e]_q$, again all of dimension $d+e-1$. The two counts
agree by
\[
  [d+e]_q \;=\; [d]_q + q^{d}\,[e]_q . \qedhere
\]
\end{proof}

Verified by enumeration in all $50$ cases with $q \in \{2,3\}$ and position
dimensions up to $4$; in the same $50$ cases the untwisted rule
$\partial(C\times D) = \partial C \times D + C \times \partial D$ \emph{fails},
so the shift is doing real work.

\begin{remark}[Why the twist]
Differentiating the second factor puts the hole in $D$; the $C$ half of the
structure has then already absorbed the extra dimension, and $\sigma C$ records
exactly that. The asymmetry of the $q$-Leibniz rule is a statement about
\emph{where the hole is}, which is invisible once one has passed to polynomials.
This is the programme working as intended.
\end{remark}

\subsection{Targets}

\begin{center}
\small
\setlength{\tabcolsep}{4pt}
\begin{tabular}{@{}>{\raggedright\arraybackslash}p{0.24\textwidth}>{\raggedright\arraybackslash}p{0.42\textwidth}>{\raggedright\arraybackslash}p{0.26\textwidth}@{}}
\toprule
identity & expected container statement & status \\
\midrule
$q$-Leibniz & Theorem~\ref{thm:leibniz} & \OK{} proved, verified \\
$q$-Pascal & split the Grassmannian container on a vector & known geometrically \\
$q$-binomial theorem & a decomposition of $\sigma^{n}$ & open \\
chain rule $\partial(C \circ D)$ & no clean $q$-chain rule exists & open \\
Pieri, Littlewood--Richardson & Hall polynomials, \S\ref{sec:hall} & open \\
Macdonald $P_\lambda(x;q,t)$ & needs a second parameter & \S\ref{sec:ceiling} \\
\bottomrule
\end{tabular}
\end{center}

The chain rule deserves emphasis: containers \emph{do} satisfy
$\partial(C \circ D) = (\partial C \circ D)\times \partial D$, and there is no clean
$q$-analogue of the chain rule for $D_q$. Either container composition deforms so as
to repair it, or the failure has a structural cause. Both answers say something the
polynomial picture cannot.

\subsection{The two-parameter ceiling}\label{sec:ceiling}

\begin{warning}
Macdonald's parameters are $(q,t)$ and \emph{neither is the field size}. In the
dictionary of \S\ref{sec:hall}, $t_{\mathrm{Mac}} = q_{\mathrm{field}}^{-1}$ and
$q_{\mathrm{Mac}} = 0$, since Hall--Littlewood is the $q_{\mathrm{Mac}} = 0$
specialisation. We have one parameter, and it is the Hecke one.
\end{warning}

The ladder of \S\ref{sec:hall} suggests the rule: \emph{one parameter per
deformation direction in the base category}. $\Vect_{\FF_q}$ has none and gives
divided powers; finite $\mathcal{O}$-modules have one nilpotent direction and give
Hall--Littlewood. Two parameters require a second direction, and there are three
known ways to supply one.

\begin{description}[leftmargin=1.6em,style=nextline,nosep]
\item[Two nilpotent directions.] Positions in finite-length $\FF_q[x,y]$-modules ---
the most container-native reading, ``positions in a two-dimensional array with decay
in each direction''. This is Haiman's setting, via the punctual Hilbert scheme. The
obstruction is sharp: $\FF_q[x,y]$-mod has global dimension $2$, so it is not
hereditary; the Ringel--Hall product is still associative but Green's theorem fails,
and polynomiality of the structure constants has no reason to survive. This is
exactly why Haiman needed geometry rather than Hall-algebra combinatorics.
\item[An elliptic curve.] Positions in $\mathrm{Coh}(E)$ for $E/\FF_q$ elliptic. The
Hall algebra is Burban--Schiffmann's elliptic Hall algebra, identified by
Schiffmann--Vasserot with the spherical double affine Hecke algebra, whose
eigenbasis on symmetric functions is Macdonald's. The two parameters are the two
Weil numbers $\alpha, \bar{\alpha}$ of the curve, with $\alpha\bar\alpha = q$. This
route works; container theory contributes only packaging.
\item[Decoupling the two ends.] $\ext{S \con P}V = \sum_s \Hom(P s, V)$ has
\emph{two} argument places and we have deformed one. Every parameter found so far
comes from $P$; the value slot $V$ has been left undeformed throughout. In the
Macdonald operator
$D_n = \sum_i \prod_{j \neq i} \frac{t x_i - x_j}{x_i - x_j}\,T_{q, x_i}$ the
parameter $t$ sits in the coefficients (the finite Hecke parameter, from the
positions) and $q$ sits in the shift operator $T_q$ --- which is
Definition~\ref{def:shift}. That is the same split as in the double affine Hecke
algebra, where $t$ is the finite Weyl parameter and $q$ the lattice-translation
parameter. Currently $x = q^{\dim V}$ ties the shift to the field size; decoupling
them is the open design question.
\end{description}

\noindent
Honest assessment: the third route is the only one in which container theory does
work rather than wrapping, and it is the least likely to be straightforward.
Confidence that decoupling the value slot yields two parameters at all: perhaps
$35\%$; that they are Macdonald's rather than an unrelated pair: lower.

\section{Status of every claim}\label{sec:status}

\begin{center}
\small
\setlength{\tabcolsep}{2.5pt}
\begin{tabular}{clcl}
\toprule
\# & claim & prior & outcome \\
\midrule
A & representation theorem (Thm.~\ref{thm:A}) & $85\%$ & \OK{} --- proved, plus 12 enumerations \\
B & fails when shapes linearise (Thm.~\ref{thm:B}) & $85\%$ & \OK{} --- proved; stronger than stated \\
D & derivative is Jackson $D_q$ (Thm.~\ref{thm:D}) & --- & \OK{} --- proved and verified \\
-- & counting series (Prop.~\ref{prop:count}) & --- & \OK{} --- verified by enumeration \\
C & two derivatives over $\FF_q$ & $60\%$ & \NO{} --- $\Vect_{\FF_q}$ is self-dual \\
C$'$ & two derivatives over non-self-dual base & new & open; hand-checked on $\vec{A_2}$ \\
H & Hall product over $\Vect_{\FF_q}$ (Prop.~\ref{prop:thin}) & --- & classical \\
HL & \emph{Day convolution $=$ Hall product} & --- & \NO{} --- \S\ref{sec:diverge}: counts splittings \\
HL$'$ & two composition laws, one Hall--Littlewood & new & Thm.~\ref{thm:HL} classical; unverified \\[-1pt]
S & filtration depth $\perp$ self-duality failure & new & Table~\ref{tab:threebases}; middle row open \\
L & adjunctions lift to containers (Thm.~\ref{thm:lift}) & new & proved; variance flips \\
I & $\Ind_* \dashv \Res_* \dashv \Ind_*$ (Cor.~\ref{cor:indres}) & new & proved from Thm.~\ref{thm:lift} \\
P & Day $=$ Levi, Hall $=$ parabolic (\S\ref{sec:indres}) & new & index computation; classical inputs \\
K & gap between $\Cont_q$ and Kuhn's $\mathcal{F}(q)$ & --- & open \\
M & Mackey as a distributive law (Q.~\ref{q:mackey}) & new & open \\
LZ & categorified $q$-Leibniz (Thm.~\ref{thm:leibniz}) & new & \OK{} proved; 50 cases verified \\
G & lifting Macdonald theory (\S\ref{sec:programme}) & new & programme; one parameter so far \\
\bottomrule
\end{tabular}
\end{center}

\section{Related work}

\begin{description}[leftmargin=1.6em,style=nextline,nosep]
\item[Abbott, Altenkirch, Ghani.] \emph{Categories of Containers} (FoSSaCS 2003);
\emph{Containers: Constructing Strictly Positive Types} (TCS 2005). The theory
being deformed.
\item[Gambino, Kock.] \emph{Polynomial functors and polynomial monads}. Source of
the local cartesian closure requirement, hence of \S\ref{sec:obstruction}.
\item[Kock.] \emph{Data types with symmetries and polynomial functors over
groupoids}. The species/quotient-container bridge; the structural model to
imitate.
\item[Kuhn.] Generic representation theory of $\mathcal{F}(q)$, the category of
functors $\finVect_{\FF_q} \to \Vect_{\FF_q}$. By Theorem~\ref{thm:B} this is the
$\Vect$-valued setting, in which container structure collapses; determining what
$\Cont_q$ looks like inside the $\Set$-valued analogue is open.
\item[Zelevinsky.] $\bigoplus_n \Rep\bigl(\GL_n(\FF_q)\bigr)$ as a positive
self-adjoint Hopf algebra, degenerating to $\bigoplus_n \Rep(S_n)$ --- symmetric
functions, hence species --- at $q=1$. The representation-theoretic target.
\item[Friedlander, Suslin.] Strict polynomial functors: the \emph{other}
linearisation, with positions a set and values a vector space.
\item[Ringel, Hall.] The algebra of Definition~\ref{def:hall}; noncommutative
exactly when the base has non-split extensions and no self-duality.
\item[Kapranov, Vasserot; Schiffmann.] The Hall product as convolution on the
stack of objects, along the span of short exact sequences --- the construction
that \S\ref{sec:diverge} says is needed in place of Day convolution.
\item[Macdonald.] \emph{Symmetric Functions and Hall Polynomials}, II--III: Hall
polynomials, the isomorphism $H(\mathcal{A}_{\mathcal{O}}) \cong \Lambda$, and
Hall--Littlewood $P_\lambda(x;t)$. The source for Theorem~\ref{thm:HL}.
\item[Green.] The unipotent character values $Q^{\lambda}_{\mu}(q)$ of
$\GL_n(\FF_q)$ as the transition matrix from power sums to Hall--Littlewood,
degenerating to the character table of $S_n$ at $q = 1$.
\item[Langer.] PhD thesis (\texttt{\char`~/git/lean/thesis-tex}): Macdonald-level
generalisation of the Kawanaka identity, with the Hall--Littlewood case $q=0$ as
precedent. The nearest local expertise on the middle row of
Table~\ref{tab:threebases}.
\end{description}

\section{Next}

\begin{enumerate}[nosep]
  \item \textbf{Question~\ref{q:closed}:} is the class of $q$-containers closed
        under Hall convolution? This decides whether \S\ref{sec:hall} is a
        statement about containers or only about $q$-species.
  \item \textbf{$\mathcal{O}$-position containers.} Enumerate both composition
        laws for $\FF_2[t]/(t^{k})$-modules, confirm Example~\ref{ex:separate},
        and check the Hall polynomials against Theorem~\ref{thm:HL}.
  \item \textbf{Settle Conjecture~\ref{conj:repaired}} by implementing
        $\Rep_{\FF_q}(\vec{A_n})$, and confirm the middle row of
        Table~\ref{tab:threebases}, isolating failure of self-duality as the sole
        obstruction to $\partial_{\mathrm{sub}} \cong \partial_{\mathrm{quot}}$.
  \item Check that Theorem~\ref{thm:A} is stable under change of base.
  \item Determine whether $\Cont_q$ sits inside a $\Set$-valued analogue of
        Kuhn's $\mathcal{F}(q)$, and measure the gap.
  \item Ask whether a nondegenerate form on the position objects makes a
        $q$-container self-dual, and whether FFD Theorem~14.4 then classifies such
        structures into three families.
\end{enumerate}

\appendix
\section{Reproducing the computations}

\begin{verbatim}
  q-bool/code/qcontainer.py   definitions, enumeration, both checkers
  q-bool/code/experiments.py  the six experiments reported above
  q-bool/code/FINDINGS.md     raw results

  $ cd q-bool/code && python3 experiments.py
\end{verbatim}

\noindent
All arithmetic is exact over prime fields; there is no sampling and no floating
point. The exhaustive checker enumerates every function family and filters by
naturality, making no use of the Yoneda lemma; the generic checker enumerates
candidates by their value on $(s,\mathrm{id}_{P s})$, verifies naturality of each
in full, and separately confirms that every element of the source is
$\ext{C}(\varphi)$ applied to a generic.

\end{document}
